\begin{abstract}
These are the expanded lecture notes from the author's mini-course during the graduate summer school of the Park City Math Institute in 2015. The main topics covered are: geometry of Springer fibers, affine Springer fibers and Hitchin fibers; representations of (affine) Weyl groups arising from these objects; relation between affine Springer fibers and orbital integrals.
\end{abstract}

%%%%%%%%%%%%%%%%%%%%%%%%%%%%%%%%%%%%%%%%%%%%%%%%%%%%%%%%%%%%%%%%%%%%%
%    
%    Make the title page
%    
%%%%%%%%%%%%%%%%%%%%%%%%%%%%%%%%%%%%%%%%%%%%%%%%%%%%%%%%%%%%%%%%%%%%%
\maketitle

\tableofcontents

%%%%%%%%%%%%%%%%%%%%%%%%%%%%%%%%%%%%%%%%%%%%%%%%%%%%%%%%%%%%%%%%%%%%%
%    
%    Your lecture notes replace the remainder of this document.
%    
%%%%%%%%%%%%%%%%%%%%%%%%%%%%%%%%%%%%%%%%%%%%%%%%%%%%%%%%%%%%%%%%%%%%%

% Authors: insert the body of your lectures here
\section{Introduction}

\subsection{Topics of these notes} These are the lectures notes from a mini-course that I gave at the PCMI graduate summer school in 2015. The goal is twofold. First I would like to introduce to the audience some interesting geometric objects that have representation-theoretic applications, and Springer fibers and their generalizations are nice examples of such. Second is to introduce orbital integrals with emphasis on its relationship with affine Springer fibers, and thereby supplying background material for B-C. Ng\^o's mini-course on the Fundamental Lemmas.

The geometric part of these lectures (everything except \S\ref{s:orb}) consists of the study of three types of ``fibers'': Springer fibers, affine Springer fibers and Hitchin fibers, with increasing complexity. We will study their geometric properties such as connectivity and irreducible components. We will construct certain group actions on these varieties, and use the action to study several nontrivial examples. Most importantly we will study certain Weyl group actions on the cohomology of these fibers which do not come from actions on the varieties themselves. The representation-theoretic significance of these three types of fibers and the analogy between them can be summarized in the following table.

\begin{table*}
	\centering
		\begin{tabular}{|c|c|c|c|}
		\hline
		& Springer fibers & affine  Springer fibers & Hitchin fibers\\
		\hline
		field & $k$ & local field $F=k\lr{t}$ & global field $k(X)$\\
		\hline
		symmetry & $W$ & $\tilW$ &  $\tilW$\\
		\hline
		extended sym& graded AHA & graded DAHA & graded DAHA\\
		\hline
		rep theory & characters of  & orbital integrals &  trace formula \\
		when $k=\FF_{q}$ & $G(k)$ & for $G(F)$ & for $G$ over $k(X)$\\
		\hline
		\end{tabular}
\end{table*}
Here $AHA$ stands for the Affine Hecke Algebra, while DAHA stands for the Double Affine Hecke Algebra; $X$ denotes an algebraic curve over $k$; $W$ and $\tilW$ are the Weyl group and extended affine Weyl group. 

In these lecture notes we do not try to give complete proofs to all statements but instead to point out interesting phenomena and examples. We do, however, give more or less complete proofs of several key results, such as 
\begin{itemize}
\item Theorem \ref{th:Spr corr} (the Springer correspondence);
\item Theorem \ref{th:ft quot} (finiteness properties of affine Springer fibers);
\item Theorem \ref{th:orb coho} (cohomological interpretation of stable orbital integrals).
\end{itemize}


\subsection{What we assume from the readers} The target readers for these lectures are beginning graduate students interested in geometric representation theory. We assume some basic algebraic geometry (scheme theory, moduli problems, point counting over a finite field, etc) though occasionally we will use the language of algebraic stacks. We also assume some Lie theory (reductive groups over an algebraic closed field and over a local field), but knowing $\GL_{n}$ and $\SL_{n}$ should be enough to understand most of these notes. 

The next remark is about the cohomology theory we use in these notes. Since we work with algebraic varieties over a general field $k$ instead of $\CC$, we will be using the \'etale cohomology with coefficients in $\ell$-adic sheaves (usually the constant sheaf $\Ql$) on these varieties. We denote the \'etale cohomology of a scheme $X$ over $k$ with $\Ql$-coefficients simply by $\cohog{*}{X}$. Readers not familiar with \'etale cohomology are encouraged to specialize to the case $k=\CC$ and understand $\cohog{*}{X}$ as the singular cohomology of $X(\CC)$ with $\QQ$-coefficients. Perverse sheaves will be used only in \S\ref{ss:Spr corr}.


\subsection*{Acknowledgement} I would like to thank the co-organizers, lecturers and the staff of the PCMI summer program in 2015. I would also like to thank the audience of my lectures for their feedback. I am especially grateful to Jingren Chi who carefully read through the first draft of these notes and provided helpful suggestions.

